% ------------------------------------------------------------------------
%  Linear solver comparison (bench_sweep.py sweep A).
%  Generated by scripts/plot_solver_bench.py -- regenerate rather than edit.
%  Include with % ------------------------------------------------------------------------
%  Linear solver comparison (bench_sweep.py sweep A).
%  Generated by scripts/plot_solver_bench.py -- regenerate rather than edit.
%  Include with % ------------------------------------------------------------------------
%  Linear solver comparison (bench_sweep.py sweep A).
%  Generated by scripts/plot_solver_bench.py -- regenerate rather than edit.
%  Include with % ------------------------------------------------------------------------
%  Linear solver comparison (bench_sweep.py sweep A).
%  Generated by scripts/plot_solver_bench.py -- regenerate rather than edit.
%  Include with \input{figures/solver_bench}
%  Needs \usepackage{pgfplots} in the preamble.
% ------------------------------------------------------------------------
\input{figures/data/solver_bench_ref}

\pgfplotsset{compat=1.16}

% Keep only the rows whose column #1 equals #2.
\pgfplotsset{
  discard if not/.style 2 args={
    x filter/.code={%
      \edef\tempa{\thisrow{#1}}\edef\tempb{#2}%
      \ifx\tempa\tempb\else\def\pgfmathresult{inf}\fi
    }
  }
}

\definecolor{solverlis}{HTML}{D95F02}
\definecolor{solveramgx}{HTML}{9DC3E0}
\definecolor{solvercfg}{HTML}{2C5F8D}

\begin{figure}[!htb]
\centering
\begin{tikzpicture}
\begin{axis}[
    ybar,
    width=9.6cm, height=6.4cm,
    bar width=26pt,
    enlarge x limits=0.18,
    ymin=0, ymax=0.383123,
    ylabel={\solverMetric{} time [s]},
    ylabel style={font=\footnotesize},
    xtick={0,1,2},
    xticklabels={{LIS},{AMGX\\(default)},{AMGX\\(tuned)}},
    xmin=-0.7, xmax=2.7,
    xticklabel style={font=\footnotesize, align=center},
    tick label style={font=\scriptsize},
    ymajorgrids=true,
    major grid style={black!12},
    % Run-to-run spread is only a few percent, so the upper whisker is 1-2 mm and
    % the lower half is buried in the bar fill. Caps nearly as wide as the bar
    % are what make it read as an error bar rather than a stray tick.
    % T-caps. rotate=90 is required: pgfplots draws the cap perpendicular to the
    % error direction by itself, but any mark size given here clobbers that, and
    % the cap then renders vertically -- which reads as a longer whisker rather
    % than a wider cap. Setting error mark explicitly does not help; the rotation
    % is what matters.
    error bars/error bar style={black!85, line width=0.8pt},
    error bars/error mark options={
      mark size=6pt, rotate=90, line width=0.8pt, black!85,
    },
    every node near coord/.append style={
      font=\tiny, color=black!60, anchor=south, yshift=5pt,
    },
]
  \addplot+[
    ybar, bar shift=0pt, fill=solverlis, draw=solverlis!70!black, line width=0.4pt,
    nodes near coords, point meta=explicit symbolic,
    error bars/.cd, y dir=both, y explicit,
  ] table[x=idx, y=value, y error minus=em, y error plus=ep, meta=iters,
          discard if not={config}{lis}] {figures/data/solver_bench.dat};
  \addplot+[
    ybar, bar shift=0pt, fill=solveramgx, draw=solveramgx!70!black, line width=0.4pt,
    nodes near coords, point meta=explicit symbolic,
    error bars/.cd, y dir=both, y explicit,
  ] table[x=idx, y=value, y error minus=em, y error plus=ep, meta=iters,
          discard if not={config}{amgx}] {figures/data/solver_bench.dat};
  \addplot+[
    ybar, bar shift=0pt, fill=solvercfg, draw=solvercfg!70!black, line width=0.4pt,
    nodes near coords, point meta=explicit symbolic,
    error bars/.cd, y dir=both, y explicit,
  ] table[x=idx, y=value, y error minus=em, y error plus=ep, meta=iters,
          discard if not={config}{amgxcfg}] {figures/data/solver_bench.dat};

\end{axis}
\end{tikzpicture}

\caption[Linear solver comparison]{\textbf{Linear solver comparison on
    \solverMolecules{}.} Median \solverMetric{} time over \solverRepeats{} runs
    per configuration, with whiskers spanning the fastest and slowest of those
    runs; the small figure above each column is the iteration count. All three
    configurations use the same discretisation and the same energy path
    (\texttt{energy\_method\,=\,1}), so the bars differ only in the linear
    solver. Note that they do not all converge to the same residual --
    table~\ref{tab:solver-bench} gives the residual each one actually reached,
    without which the times are not comparable.}
\label{fig:solver-bench}
\end{figure}

%  Needs \usepackage{pgfplots} in the preamble.
% ------------------------------------------------------------------------
% generated by scripts/plot_solver_bench.py -- do not edit
\def\solverMetric{solve stage}
\def\solverRepeats{3}
\def\solverUnit{seconds}
\def\solverMolecules{1CCM, 6VYB, 1vsz}


\pgfplotsset{compat=1.16}

% Keep only the rows whose column #1 equals #2.
\pgfplotsset{
  discard if not/.style 2 args={
    x filter/.code={%
      \edef\tempa{\thisrow{#1}}\edef\tempb{#2}%
      \ifx\tempa\tempb\else\def\pgfmathresult{inf}\fi
    }
  }
}

\definecolor{solverlis}{HTML}{D95F02}
\definecolor{solveramgx}{HTML}{9DC3E0}
\definecolor{solvercfg}{HTML}{2C5F8D}

\begin{figure}[!htb]
\centering
\begin{tikzpicture}
\begin{axis}[
    ybar,
    width=9.6cm, height=6.4cm,
    bar width=26pt,
    enlarge x limits=0.18,
    ymin=0, ymax=0.383123,
    ylabel={\solverMetric{} time [s]},
    ylabel style={font=\footnotesize},
    xtick={0,1,2},
    xticklabels={{LIS},{AMGX\\(default)},{AMGX\\(tuned)}},
    xmin=-0.7, xmax=2.7,
    xticklabel style={font=\footnotesize, align=center},
    tick label style={font=\scriptsize},
    ymajorgrids=true,
    major grid style={black!12},
    % Run-to-run spread is only a few percent, so the upper whisker is 1-2 mm and
    % the lower half is buried in the bar fill. Caps nearly as wide as the bar
    % are what make it read as an error bar rather than a stray tick.
    % T-caps. rotate=90 is required: pgfplots draws the cap perpendicular to the
    % error direction by itself, but any mark size given here clobbers that, and
    % the cap then renders vertically -- which reads as a longer whisker rather
    % than a wider cap. Setting error mark explicitly does not help; the rotation
    % is what matters.
    error bars/error bar style={black!85, line width=0.8pt},
    error bars/error mark options={
      mark size=6pt, rotate=90, line width=0.8pt, black!85,
    },
    every node near coord/.append style={
      font=\tiny, color=black!60, anchor=south, yshift=5pt,
    },
]
  \addplot+[
    ybar, bar shift=0pt, fill=solverlis, draw=solverlis!70!black, line width=0.4pt,
    nodes near coords, point meta=explicit symbolic,
    error bars/.cd, y dir=both, y explicit,
  ] table[x=idx, y=value, y error minus=em, y error plus=ep, meta=iters,
          discard if not={config}{lis}] {figures/data/solver_bench.dat};
  \addplot+[
    ybar, bar shift=0pt, fill=solveramgx, draw=solveramgx!70!black, line width=0.4pt,
    nodes near coords, point meta=explicit symbolic,
    error bars/.cd, y dir=both, y explicit,
  ] table[x=idx, y=value, y error minus=em, y error plus=ep, meta=iters,
          discard if not={config}{amgx}] {figures/data/solver_bench.dat};
  \addplot+[
    ybar, bar shift=0pt, fill=solvercfg, draw=solvercfg!70!black, line width=0.4pt,
    nodes near coords, point meta=explicit symbolic,
    error bars/.cd, y dir=both, y explicit,
  ] table[x=idx, y=value, y error minus=em, y error plus=ep, meta=iters,
          discard if not={config}{amgxcfg}] {figures/data/solver_bench.dat};

\end{axis}
\end{tikzpicture}

\caption[Linear solver comparison]{\textbf{Linear solver comparison on
    \solverMolecules{}.} Median \solverMetric{} time over \solverRepeats{} runs
    per configuration, with whiskers spanning the fastest and slowest of those
    runs; the small figure above each column is the iteration count. All three
    configurations use the same discretisation and the same energy path
    (\texttt{energy\_method\,=\,1}), so the bars differ only in the linear
    solver. Note that they do not all converge to the same residual --
    table~\ref{tab:solver-bench} gives the residual each one actually reached,
    without which the times are not comparable.}
\label{fig:solver-bench}
\end{figure}

%  Needs \usepackage{pgfplots} in the preamble.
% ------------------------------------------------------------------------
% generated by scripts/plot_solver_bench.py -- do not edit
\def\solverMetric{solve stage}
\def\solverRepeats{3}
\def\solverUnit{seconds}
\def\solverMolecules{1CCM, 6VYB, 1vsz}


\pgfplotsset{compat=1.16}

% Keep only the rows whose column #1 equals #2.
\pgfplotsset{
  discard if not/.style 2 args={
    x filter/.code={%
      \edef\tempa{\thisrow{#1}}\edef\tempb{#2}%
      \ifx\tempa\tempb\else\def\pgfmathresult{inf}\fi
    }
  }
}

\definecolor{solverlis}{HTML}{D95F02}
\definecolor{solveramgx}{HTML}{9DC3E0}
\definecolor{solvercfg}{HTML}{2C5F8D}

\begin{figure}[!htb]
\centering
\begin{tikzpicture}
\begin{axis}[
    ybar,
    width=9.6cm, height=6.4cm,
    bar width=26pt,
    enlarge x limits=0.18,
    ymin=0, ymax=0.383123,
    ylabel={\solverMetric{} time [s]},
    ylabel style={font=\footnotesize},
    xtick={0,1,2},
    xticklabels={{LIS},{AMGX\\(default)},{AMGX\\(tuned)}},
    xmin=-0.7, xmax=2.7,
    xticklabel style={font=\footnotesize, align=center},
    tick label style={font=\scriptsize},
    ymajorgrids=true,
    major grid style={black!12},
    % Run-to-run spread is only a few percent, so the upper whisker is 1-2 mm and
    % the lower half is buried in the bar fill. Caps nearly as wide as the bar
    % are what make it read as an error bar rather than a stray tick.
    % T-caps. rotate=90 is required: pgfplots draws the cap perpendicular to the
    % error direction by itself, but any mark size given here clobbers that, and
    % the cap then renders vertically -- which reads as a longer whisker rather
    % than a wider cap. Setting error mark explicitly does not help; the rotation
    % is what matters.
    error bars/error bar style={black!85, line width=0.8pt},
    error bars/error mark options={
      mark size=6pt, rotate=90, line width=0.8pt, black!85,
    },
    every node near coord/.append style={
      font=\tiny, color=black!60, anchor=south, yshift=5pt,
    },
]
  \addplot+[
    ybar, bar shift=0pt, fill=solverlis, draw=solverlis!70!black, line width=0.4pt,
    nodes near coords, point meta=explicit symbolic,
    error bars/.cd, y dir=both, y explicit,
  ] table[x=idx, y=value, y error minus=em, y error plus=ep, meta=iters,
          discard if not={config}{lis}] {figures/data/solver_bench.dat};
  \addplot+[
    ybar, bar shift=0pt, fill=solveramgx, draw=solveramgx!70!black, line width=0.4pt,
    nodes near coords, point meta=explicit symbolic,
    error bars/.cd, y dir=both, y explicit,
  ] table[x=idx, y=value, y error minus=em, y error plus=ep, meta=iters,
          discard if not={config}{amgx}] {figures/data/solver_bench.dat};
  \addplot+[
    ybar, bar shift=0pt, fill=solvercfg, draw=solvercfg!70!black, line width=0.4pt,
    nodes near coords, point meta=explicit symbolic,
    error bars/.cd, y dir=both, y explicit,
  ] table[x=idx, y=value, y error minus=em, y error plus=ep, meta=iters,
          discard if not={config}{amgxcfg}] {figures/data/solver_bench.dat};

\end{axis}
\end{tikzpicture}

\caption[Linear solver comparison]{\textbf{Linear solver comparison on
    \solverMolecules{}.} Median \solverMetric{} time over \solverRepeats{} runs
    per configuration, with whiskers spanning the fastest and slowest of those
    runs; the small figure above each column is the iteration count. All three
    configurations use the same discretisation and the same energy path
    (\texttt{energy\_method\,=\,1}), so the bars differ only in the linear
    solver. Note that they do not all converge to the same residual --
    table~\ref{tab:solver-bench} gives the residual each one actually reached,
    without which the times are not comparable.}
\label{fig:solver-bench}
\end{figure}

%  Needs \usepackage{pgfplots} in the preamble.
% ------------------------------------------------------------------------
% generated by scripts/plot_solver_bench.py -- do not edit
\def\solverMetric{solve stage}
\def\solverRepeats{3}
\def\solverUnit{seconds}
\def\solverMolecules{1CCM, 6VYB, 1vsz}


\pgfplotsset{compat=1.16}

% Keep only the rows whose column #1 equals #2.
\pgfplotsset{
  discard if not/.style 2 args={
    x filter/.code={%
      \edef\tempa{\thisrow{#1}}\edef\tempb{#2}%
      \ifx\tempa\tempb\else\def\pgfmathresult{inf}\fi
    }
  }
}

\definecolor{solverlis}{HTML}{D95F02}
\definecolor{solveramgx}{HTML}{9DC3E0}
\definecolor{solvercfg}{HTML}{2C5F8D}

\begin{figure}[!htb]
\centering
\begin{tikzpicture}
\begin{axis}[
    ybar,
    width=9.6cm, height=6.4cm,
    bar width=26pt,
    enlarge x limits=0.18,
    ymin=0, ymax=0.383123,
    ylabel={\solverMetric{} time [s]},
    ylabel style={font=\footnotesize},
    xtick={0,1,2},
    xticklabels={{LIS},{AMGX\\(default)},{AMGX\\(tuned)}},
    xmin=-0.7, xmax=2.7,
    xticklabel style={font=\footnotesize, align=center},
    tick label style={font=\scriptsize},
    ymajorgrids=true,
    major grid style={black!12},
    % Run-to-run spread is only a few percent, so the upper whisker is 1-2 mm and
    % the lower half is buried in the bar fill. Caps nearly as wide as the bar
    % are what make it read as an error bar rather than a stray tick.
    % T-caps. rotate=90 is required: pgfplots draws the cap perpendicular to the
    % error direction by itself, but any mark size given here clobbers that, and
    % the cap then renders vertically -- which reads as a longer whisker rather
    % than a wider cap. Setting error mark explicitly does not help; the rotation
    % is what matters.
    error bars/error bar style={black!85, line width=0.8pt},
    error bars/error mark options={
      mark size=6pt, rotate=90, line width=0.8pt, black!85,
    },
    every node near coord/.append style={
      font=\tiny, color=black!60, anchor=south, yshift=5pt,
    },
]
  \addplot+[
    ybar, bar shift=0pt, fill=solverlis, draw=solverlis!70!black, line width=0.4pt,
    nodes near coords, point meta=explicit symbolic,
    error bars/.cd, y dir=both, y explicit,
  ] table[x=idx, y=value, y error minus=em, y error plus=ep, meta=iters,
          discard if not={config}{lis}] {figures/data/solver_bench.dat};
  \addplot+[
    ybar, bar shift=0pt, fill=solveramgx, draw=solveramgx!70!black, line width=0.4pt,
    nodes near coords, point meta=explicit symbolic,
    error bars/.cd, y dir=both, y explicit,
  ] table[x=idx, y=value, y error minus=em, y error plus=ep, meta=iters,
          discard if not={config}{amgx}] {figures/data/solver_bench.dat};
  \addplot+[
    ybar, bar shift=0pt, fill=solvercfg, draw=solvercfg!70!black, line width=0.4pt,
    nodes near coords, point meta=explicit symbolic,
    error bars/.cd, y dir=both, y explicit,
  ] table[x=idx, y=value, y error minus=em, y error plus=ep, meta=iters,
          discard if not={config}{amgxcfg}] {figures/data/solver_bench.dat};

\end{axis}
\end{tikzpicture}

\caption[Linear solver comparison]{\textbf{Linear solver comparison on
    \solverMolecules{}.} Median \solverMetric{} time over \solverRepeats{} runs
    per configuration, with whiskers spanning the fastest and slowest of those
    runs; the small figure above each column is the iteration count. All three
    configurations use the same discretisation and the same energy path
    (\texttt{energy\_method\,=\,1}), so the bars differ only in the linear
    solver. Note that they do not all converge to the same residual --
    table~\ref{tab:solver-bench} gives the residual each one actually reached,
    without which the times are not comparable.}
\label{fig:solver-bench}
\end{figure}
